\chapter[The coherence engine]{The coherence engine: CSV finite-state machines}\label{ch:coherence-engine}

\whatyoufind{Why a coherence protocol in Octopus is a spreadsheet and not a program; the exact
format of that spreadsheet and the one rule about it that silently breaks protocols; the
vocabulary of events a controller can see and actions it can perform; the families that ship
and how a snooping protocol differs from a directory one; and a worked example, MI, added
without a line of C++.}

\section{Coherence as data}

A cache controller's coherence logic is a finite-state machine read at start-up from a CSV
file. Rows are states, stable and transient; columns are coherence events; each cell lists the
actions and the next state for that (state, event) pair. Adding or modifying a protocol means
editing the table --- no C++ change and no recompilation. The same engine, \code{FSMReader} plus
a per-family protocol handler, drives MSI, MESI and MOESI in their snooping and directory forms,
and the predictable protocols built on Octopus (\cref{ch:introduction}).

The split pays off in code size. Because the handler only translates a small set of protocol
actions into controller actions, a protocol that extends a family it already knows costs little:
\cref{tab:loc} contrasts the lines of code per protocol with gem5's Ruby.

\begin{table}[htbp]\centering\small
\caption{Lines of code per coherence protocol, Octopus versus gem5's Ruby. In Octopus the count
is the CSV table plus the handler code specific to that protocol; MESI and MOESI extend MSI's
handler.}\label{tab:loc}
\begin{tabular}{@{}lrrrr@{}}
\toprule & MSI & MESI & MOESI & PENDULUM \\ \midrule
gem5 (Ruby) & 800 & 1400 & 1300 & 2000 \\
\textbf{Octopus} & \textbf{300} & \textbf{70} & \textbf{120} & \textbf{500} \\
\bottomrule
\end{tabular}
\end{table}

\section{The table}\label{sec:fsm-format}

\begin{figure}[htbp]\centering
\begin{adjustbox}{max width=\textwidth}% An excerpt of MESI_LLC.csv with the four things a reader must know about the format.
\begin{tikzpicture}[font=\sffamily\footnotesize, >=Latex,
  cell/.style={minimum height=6mm, inner sep=3pt, anchor=center, font=\ttfamily\footnotesize},
  hdr/.style={cell, fill=octmem, font=\sffamily\bfseries\footnotesize},
  st/.style={cell, fill=octctl, font=\ttfamily\bfseries\footnotesize},
  note/.style={draw=octquiet, fill=octpaper, rounded corners=2pt, inner sep=4pt, align=left, font=\sffamily\scriptsize, text width=62mm},
  pin/.style={->, octquiet, thick}]
  \matrix (m) [matrix of nodes, nodes={cell}, column sep=-\pgflinewidth, row sep=-\pgflinewidth, nodes in empty cells,
    column 1/.style={nodes={st}}, row 1/.style={nodes={hdr}}] {
    State & stable & GetS & GetM & Replacement & Own\_Invalidation \\
    N     & 1 & GetData/SetOwner/NE\_d & GetData/SetOwner/NM\_d & & SaveData/N \\
    NE\_d & 0 & Stall/ & Stall/ & Stall/ & Fault/ \\
    I     & 1 & SendExeclusiveData/SetOwner/EorM & SendData/SetOwner/EorM & IssueInv/ & WriteBack/N \\
    EorM  & 1 & ClearOwner/S\_d & SetOwner/ & IssueInv/ & ClearOwner/MN\_d \\
  };
  \draw[octrule] (m-1-1.south west) -- (m-1-6.south east);
  \draw[octrule] (m-1-1.north west) rectangle (m-5-6.south east);
  % callouts above
  \node[note, above=10mm of m-1-4.north, anchor=south] (nev) {\textbf{a column is an event}, one of the \code{EventNum} list at the top of the file; the handler's \code{readEvent} maps each message to one of them};
  \draw[pin] (nev.south) -- (m-1-4.north);
  \node[note, above=10mm of m-1-1.north, anchor=south west, xshift=-4mm] (nst) {\textbf{a row is a state}, stable (1) or transient (0); its \emph{position} in this table must equal its \code{StateNum}, so a new state's row goes \textbf{last}};
  \draw[pin] (nst.south) -- (m-1-1.north);
  % callouts below
  \node[note, below=10mm of m-5-3.south, anchor=north] (ncell) {\textbf{a cell is \code{actions/next}}: do \code{SendData}, then \code{SetOwner}, then go to \code{EorM}. No trailing state (\code{IssueInv/}) means \emph{stay}; an empty cell \emph{ignores} the event};
  \draw[pin] (ncell.north) -- (m-4-4.south);
  \node[note, below=10mm of m-5-6.south, anchor=north east, xshift=4mm] (nf) {\textbf{\code{Fault/}} marks a pair that must be unreachable; the run halts and names it, which is how an incomplete table fails loudly};
  \draw[pin] (nf.north) -- (m-3-6.south);
\end{tikzpicture}
\end{adjustbox}
\caption{An excerpt of \code{MESI\_LLC.csv} with the four things to know about the format.
\Cref{app:fsm-tables} prints every shipped table in full.}
\label{fig:fsm-annotated}
\end{figure}

A file in \file{Protocols_FSM/} has four blocks (\cref{fig:fsm-annotated}):

\begin{enumerate}
\item \textbf{Events}, numbered from 0: the columns of the transition table. For the snooping
  L1 they are exactly \evt{Load}, \evt{Store}, \evt{Replacement}, \evt{Own\_GetS},
  \evt{Own\_GetM}, \evt{Own\_PutM}, \evt{Other\_GetS}, \evt{Other\_GetM}, \evt{Other\_PutM},
  \evt{OwnData}, \evt{Invalidation} and \evt{OwnData\_Execlusive}; for the snooping LLC
  \evt{GetS}, \evt{GetM}, \evt{Replacement}, \evt{PutM\_fromOwner}, \evt{PutM\_fromNonOwner},
  \evt{Data\_fromLowerInterface}, \evt{Data\_fromUpperInterface} and \evt{Own\_Invalidation}.
\item \textbf{Actions}, numbered from 0: what a cell may ask for. The L1 offers \act{Stall},
  \act{Hit}, \act{GetS}, \act{GetM}, \act{PutM}, \act{Data2Req}, \act{Data2Both}, \act{SaveReq},
  \act{Fault} and \act{removeSavedReq}; the LLC offers \act{Stall}, \act{GetData},
  \act{SendData}, \act{SaveData}, \act{SetOwner}, \act{ClearOwner}, \act{IssueInv},
  \act{WriteBack}, \act{Fault} and \act{SendExeclusiveData}. Each type of controller offers its
  list independently of the protocol, which is what lets one controller class run several
  protocols.
\item \textbf{States}, numbered from 0, stable and transient alike. Transient names follow the
  textbook convention: \state{IS\_ad} is ``going from I to S, awaiting the bus (a) and the data
  (d)'', \state{IM\_dI} is ``going to M, data pending, but invalidated meanwhile''.
\item \textbf{The transition table}: a header row \code{State, stable, isDataValid, <events>},
  then one row per state. \code{stable} marks the stable states; \code{isDataValid} says whether
  the line's data may be read in that state. A cell is \code{Action1/Action2/NextState}. A cell
  whose last element is an action, not a state (\code{IssueInv/}), stays in the current state;
  an empty cell ignores the event; \code{Fault/} declares the pair unreachable.
\end{enumerate}

\begin{trapbox}[Row order is the state number]
\code{FSMReader} indexes transition rows by \emph{position}, not by name. When you add a state,
its row must be appended \textbf{last}, so that its position matches the number you gave it in
the state list. Put it anywhere else and the FSM silently runs another state's transitions ---
no error, just wrong results. This is the single most likely way a protocol edit goes wrong.
\end{trapbox}

\paragraph{Reading a cell.} In \cref{fig:fsm-annotated}, state \state{I} (the LLC has the line,
no L1 owns it) on \evt{GetM} reads \code{SendData/SetOwner/EorM}: send the data to the
requester, record it as the owner, and move to \state{EorM}. The same state on
\evt{Replacement} reads \code{IssueInv/}: broadcast an invalidation and stay --- the eviction
completes later, when the LLC's own copy of that invalidation comes back as
\evt{Own\_Invalidation} and the cell \code{WriteBack/N} runs. \Cref{ch:llc-datapath} draws those
two paths.

\paragraph{Fault is an audit.} A cell of \code{Fault/} means ``this (state, event) should be
unreachable''. The handler halts the run and reports the pair rather than mis-routing silently.
That is what makes a protocol auditable: an incomplete or wrong table fails loudly on the first
trace that exercises the hole, and the Debugger's transition trace (\cref{ch:monitoring}) shows
the exact sequence that reached it.

\section{Stable and transient states}

A snooping protocol on a split-transaction bus is mostly transient states. A request is issued
(\act{GetS}), travels the request lane, is seen back by its own issuer (\evt{Own\_GetS}: the
line is now ordered on the bus), and some time later the data arrives (\evt{OwnData}); between
those moments other cores' requests to the same line are seen (\evt{Other\_GetS},
\evt{Other\_GetM}) and may have to be remembered (\act{SaveReq}) and answered when the data
finally comes (\act{Data2Req}, \act{Data2Both}). Each such interleaving is a row. The shipped
MESI L1 has 4 stable and 20 transient states; its LLC 4 stable and 7 transient
(\cref{app:fsm-tables}).

The LLC side uses states a textbook does not name. \state{N} is ``not present'';
\state{EorM} is ``an L1 owns this line, I do not know in which of E or M''; \state{MN\_d} is
``evicting, waiting for the owner's data''. They exist because the LLC is inclusive and tracks
one owner per line, and they are where most of the subtle behaviour of \cref{ch:llc-datapath}
lives.

\section{Snooping and directory}

\begin{figure}[htbp]\centering
\begin{adjustbox}{max width=\textwidth}% One GetM from core A while core B shares the line: snooping (left) and directory (right).
\begin{tikzpicture}[font=\sffamily\scriptsize, >=Latex,
  head/.style={draw=octrule, fill=octmem, rounded corners=2pt, minimum width=16mm, minimum height=6mm, font=\sffamily\footnotesize\bfseries},
  life/.style={octquiet, densely dashed},
  m/.style={->, thick},
  lbl/.style={font=\sffamily\scriptsize, fill=octpaper, inner sep=1pt},
  pill/.style={draw, fill=white, rounded corners=3pt, inner sep=2pt, font=\ttfamily\scriptsize},
  ttl/.style={font=\sffamily\small\bfseries, text=octink}]
  % ---------- snooping ----------
  \begin{scope}
    \node[ttl] at (28mm, 8mm) {Snooping (shipped MESI)};
    \foreach \x/\n/\t in {0/A/Core A L1, 19/B/Core B L1, 38/bus/bus[0], 57/llc/LLC}
      { \node[head] (s\n) at (\x mm, 0) {\t}; \draw[life] (s\n.south) -- ++(0,-46mm); }
    \draw[m, octpurple] (sA) ++(0,-8mm) -- node[lbl, above] {GetM X, request lane: broadcast} ++(57mm, 0);
    \node[pill, draw=octamber] at (19mm, -14mm) {B snoops: S $\rightarrow$ I};
    \node[pill, draw=octteal] at (57mm, -14mm) {LLC: SendData, SetOwner A};
    \draw[m, octpurple] (sllc) ++(0,-22mm) -- node[lbl, above] {X, response lane} ++(-57mm, 0);
    \node[pill, draw=octteal] at (0mm, -29mm) {A: IM\_d $\rightarrow$ M};
    \node[align=left, text=octquiet, text width=60mm] at (30mm, -40mm) {one broadcast, every sharer invalidates itself by seeing it; the order every controller sees on the bus is the order of the protocol};
  \end{scope}
  % ---------- directory ----------
  \begin{scope}[xshift=84mm]
    \node[ttl] at (28mm, 8mm) {Directory (shipped MSI)};
    \foreach \x/\n/\t in {0/A/Core A L1, 19/B/Core B L1, 38/home/LLC + directory, 57/mem/memory}
      { \node[head] (d\n) at (\x mm, 0) {\t}; \draw[life] (d\n.south) -- ++(0,-46mm); }
    \draw[m, octpurple] (dA) ++(0,-8mm) -- node[lbl, above] {GetM X, point-to-point} ++(38mm, 0);
    \node[pill, draw=octteal] at (38mm, -13mm) {sharers = \{B\}};
    \draw[m, octred] (dhome) ++(0,-18mm) -- node[lbl, above] {Inv X} ++(-19mm, 0);
    \draw[m, octpurple] (dhome) ++(0,-23mm) -- node[lbl, above] {Data X, acks = 1} ++(-38mm, 0);
    \node[pill, draw=octamber] at (19mm, -27mm) {B: S $\rightarrow$ I};
    \draw[m, octteal] (dB) ++(0,-31mm) -- node[lbl, above] {Inv-Ack} ++(-19mm, 0);
    \node[pill, draw=octteal] at (0mm, -37mm) {A: all acks in $\rightarrow$ M};
    \node[align=left, text=octquiet, text width=60mm] at (30mm, -43mm) {no broadcast: the home knows the sharers and forwards; the requester counts acks. Works on any point-to-point network};
  \end{scope}
\end{tikzpicture}
\end{adjustbox}
\caption{One \code{GetM} from core A while core B shares the line, under the shipped snooping
MESI (left) and the shipped directory MSI (right).}
\label{fig:snoop-vs-directory}
\end{figure}

The two families differ in who learns about a request (\cref{fig:snoop-vs-directory}). Under
snooping every L1 sees every request on the request lane and acts on its own copy; the order of
the bus is the order of the protocol, and the LLC's job is to supply data when no L1 owns the
line and to track the owner when one does. Under a directory the home --- the LLC-side
\code{CacheControllerDirectory} --- keeps a sharer list per line, forwards each request only to
the caches that need to act, and the requester counts the acknowledgements; nothing is
broadcast, which is what lets a directory run on a point-to-point network such as the mesh
(\cref{ch:interconnects}). The same controllers, queues and data handlers serve both; only the
protocol handler and the tables differ.

\begin{table}[htbp]\centering\small
\caption{The protocol families that ship, and the preset that selects each. The L1 and LLC
tables of a family are a matched pair.}\label{tab:families}
\footnotesize
\begin{tabularx}{\textwidth}{@{}p{0.17\textwidth}p{0.27\textwidth}p{0.2\textwidth}Y@{}}
\toprule family & L1 table & LLC table & selected by \\ \midrule
snoop MSI & \file{MSI_splitBus_snooping} & \file{MSI_LLC} & \file{MultiCoreSystem_Snoop.csv} \\
snoop MESI & \file{MESI_splitBus_snooping} & \file{MESI_LLC} & \file{MultiCoreSystem.csv} (the default) \\
snoop MOESI & \file{MOESI_splitBus_snooping} & \file{MOESI_LLC} & no preset \\
directory MSI & \file{MSI_directory} & \file{MSI_LLC_directory} & \file{MultiCoreSystem_Directory.csv} \\
directory MESI, MOESI & \file{MESI_directory}, \file{MOESI_directory} & \file{MESI_LLC_directory}, \file{MOESI_LLC_directory} & no preset \\
predictable & \file{PMSI}, \file{PMESI}, \file{PMSI_asterisk}, \file{PMESI_asterisk}, \file{Pendulum} & their \file{_LLC} tables & research configurations \\
\bottomrule
\end{tabularx}
\end{table}

\begin{trapbox}[{A protocol is five keys, and the tables are a matched pair}]
Selecting a protocol sets the controller class, the L1 and LLC \code{protocol\_type} and the L1
and LLC \code{fsm\_filename} together (\cref{ch:configuration}). A MESI L1 with an MSI LLC
faults with ``Wrong destination''; overriding some of the five from the command line faults or
crashes. Switch by preset, not by hand.
\end{trapbox}

\section{A worked example: MI without C++}

MI is the simplest coherence protocol: a line is either Invalid or Modified, so every access,
even a load, takes it exclusive. It is MSI with the \state{S} state and its transients removed
and \evt{Load} behaving like \evt{Store}: 9 states instead of 22. Because MI uses only a subset
of MSI's events and actions, it runs on the existing \code{SNOOP\_MSI} handler and the existing
\code{MSI\_LLC}: the entire protocol is one CSV, \file{Protocols_FSM/MI_splitBus_snooping.csv}.

\begin{table}[htbp]\centering\small
\caption{The two stable rows of \code{MI\_splitBus\_snooping.csv}; the columns not shown are
those of MSI.}\label{tab:mi-rows}
\begin{tabular}{@{}>{\ttfamily}l>{\ttfamily}c>{\ttfamily}c>{\ttfamily}l>{\ttfamily}l>{\ttfamily}c>{\ttfamily}l>{\ttfamily}l@{}}
\toprule
\normalfont\textbf{State} & \normalfont\textbf{stable} & \normalfont\textbf{isDataValid} & \normalfont\textbf{Load} & \normalfont\textbf{Store} & \ldots & \normalfont\textbf{OwnData} & \normalfont\textbf{Invalidation} \\ \midrule
I & 1 & 0 & GetM/IM\_ad & GetM/IM\_ad & \ldots & Fault/ & \\
M & 1 & 1 & Hit/ & Hit/ & \ldots & Fault/ & Data2Req/I \\
\bottomrule
\end{tabular}
\end{table}

Selecting it is one override on top of the MSI preset:

\begin{lstlisting}[style=shell]
./build/Octopus_Simulator -s MultiCoreSystem \
  -p "workload_path(s)=BMs/eembc-traces/mi_pingpong/" \
  -p "cache_controller[*].fsm_filename(s)=MI_splitBus_snooping"
\end{lstlisting}

On the \code{mi\_pingpong} toy workload, where cores 0 and 1 repeatedly read one line, MESI lets
the two cores share it and finishes in 332 cycles; MI makes the line ping-pong exclusively
between them and finishes in 482. The Logger shows the cost; the Debugger's address-filtered
transition trace shows the cause --- \code{I --Load--> IM\_ad} under MI against
\code{I --Load--> IS\_ad} under MESI. One new file, one override, no compiler.

\begin{wherebox}
\begin{itemize}[nosep,leftmargin=1.2em]
\item \file{Protocols_FSM/} --- the tables, one CSV per agent
\item \file{header/FSMReader.h}, \file{src/FSMReader.cpp} --- the parser and \code{getTransition}
\item \file{src/Protocols/MSIProtocol.cpp}, \file{src/Protocols/MESIProtocol.cpp}, \file{src/Protocols/MOESIProtocol.cpp} --- the snooping L1 handlers
\item \file{src/Protocols/LLCMSIProtocol.cpp}, \file{src/Protocols/LLCMESIProtocol.cpp} --- the snooping LLC handlers (MOESI's LLC runs on the MESI handler)
\item \file{src/Protocols/MSIDirectory.cpp}, \file{src/Protocols/MESIDirectory.cpp}, \file{src/Protocols/MOESIDirectory.cpp} --- the directory L1 handlers
\item \file{src/Protocols/LLCMSIDirectory.cpp}, \file{src/Protocols/LLCMESIDirectory.cpp}, \file{src/Protocols/LLCMOESIDirectory.cpp} --- the directory home handlers
\item \file{docs/AddingAProtocol.md} --- the MI example in full
\item \file{demo_protocol.sh} --- runs it
\end{itemize}
\end{wherebox}
